% ====================================================================
% FILE: content/15_falsifiability_extended.tex
% Расширяемая фальсифицируемость и экспериментальная программа
% Core 1.3.3 — canonical extended version
% ====================================================================

\section{Расширяемая фальсифицируемость и экспериментальная программа}
\label{sec:falsifiability-extended}

В этом разделе восстанавливается и расширяется фреймворк фальсифицируемости
для Онтологии континуумов (OC). Формулируются явные структурные
предсказания, критерии фальсификации и экспериментальные программы по всей
иерархии уровней \(K_2\)–\(K_{12}\), используя только аксиомы и
определения Core~1.3.3.

Цель не в составлении полного каталога экспериментов, а в том, чтобы показать,
что OC структурно фальсифицируема: она делает междоменные обязательства
о порогах, циклах, размерностных переходах и граничном поведении, которые
в принципе могут быть опровергнуты эмпирическими и теоретическими данными.

% ====================================================================
\subsection{Смысл фальсифицируемости для структурных теорий}
% ====================================================================

Для структурной теории фальсифицируемость не сводится к отдельным числовым
предсказаниям. Напротив, фальсификация происходит, когда структурные
отношения между
\[
  (\Omega, A, P(t), J(t),
   \Theta, \partial\Omega, C, k(t))
\]
нарушаются в способах, несовместимых с аксиомами и теоремами OC.

\paragraph{Структурная и феноменологическая фальсифицируемость.}
\begin{itemize}
  \item \emph{Феноменологическая фальсифицируемость} касается прямых
        несоответствий между выходами модели и наблюдениями (например,
        ошибочная диаграмма фаз, неверный масштабный показатель).
  \item \emph{Структурная фальсифицируемость} касается нарушений
        универсальных ограничений континуумов: порогов, циклов,
        размерностной монотонности, совместимости вложения, условий
        схлопывания и возрождения.
\end{itemize}
OC в первую очередь работает на структурном уровне; феноменологические модели
должны проходить через структурные требования OC, чтобы считаться корректными
инстанцированиями.

\paragraph{Каналы структурной фальсификации.}
Континуумная модель, претендующая на инстанцирование OC, может быть
фальсифицирована, если:
\begin{enumerate}
  \item
        систематически происходят нарушения порогов: наблюдаемые системы
        остаются жизнеспособными при нарушении предложенных
        \(\Theta_{\mathrm{exist}}\), \(\Theta_{\mathrm{stab}}\),
        \(\Theta_{\mathrm{crit}}\) или \(\Theta_{\mathrm{death}}\);
  \item
        исчезновение циклов противоречит жизнеспособности: системы
        поддерживают длительную организацию без какого-либо идентифицируемого
        комплекса циклов \(C\), удовлетворяющего условиям OC;
  \item
        наблюдаются запрещенные размерностные переходы: эффективная размерность
        увеличивается без новых осей в пространстве вложения или без пересечения
        размерностного порога \(\Theta_{\mathrm{dim}}\);
  \item
        требуются несовместимые кортежи: успешные модели требуют взаимно
        несовместимых версий канонического OC-состояния для одной и той же
        системы.
\end{enumerate}

\paragraph{Структурная против эмпирическая неудача.}
Эмпирическая неудача конкретной предметной модели не обязательно сразу
фальсифицирует OC; структурная теория ставится под вопрос только если
эмпирический успех альтернативных моделей требует нарушений структурных
ограничений OC. Наоборот, если в данной области повторно наблюдаются
структуры, запрещенные OC (например, устойчивые континуумы без поддерживающих
циклов), OC претерпевает прямую фальсификацию.

\paragraph{Междоменная логика фальсификации.}
Поскольку утверждения OC универсальны, структурная фальсификация в одной
области (прочно установившаяся) распространяется на все области, которые
опираются на те же аксиомы. Это делает OC более фальсифицируемой, чем
локальные теории:
\begin{itemize}
  \item если в контролируемой физической системе оказывается нарушенной
        монотонность размерности, её следует пересмотреть и для когнитивных,
        и для социальных континуумов;
  \item если жизнеспособные протоклетки демонстрируются без любой
        структуры границы, совместимой с \(\partial\Omega\), затрагивается
        вся иерархия уровней.
\end{itemize}

Остальная часть главы формулирует явные предсказания и тесты.

% ====================================================================
\subsection{Предсказания и тесты в физике
            (\texorpdfstring{$K_2$}{K\_2})}
% ====================================================================

Физические континуумы предоставляют наиболее классическую площадку тестов OC на
уровне \(K_2\).

\paragraph{Структурные предсказания.}
\begin{itemize}
  \item \textbf{P2.1 (Критические пороги).}
        Фазовые переходы соответствуют структурным пересечениям
        \(\Theta_{\mathrm{crit}}\). Ни один резкий дальнодействующий
        переход упорядоченности не может происходить без соответствующей
        пороговой поверхности в \(\Omega_2\).
  \item \textbf{P2.2 (Монотонность перколяции).}
        Для систем на основе связности с управляющим параметром \(p\),
        глобальная жизнеспособность (\(k_2>0\)) требует \(p \ge p_c\);
        при \(p<p_c\) не может существовать бесконечный кластер,
        совместимый с \(\Omega_2\).
  \item \textbf{P2.3 (Размерностная монотонность).}
        Ни один физический континуум не может повысить эффективную размерность
        без активации по меньшей мере одной новой оси в пространстве вложения
        \(M_2\) и пересечения размерностного порога
        \(\Theta_{\mathrm{dim}}\).
  \item \textbf{P2.4 (Схлопывание через границу).}
        Схлопывание упорядоченных фаз происходит через разрушение поддерживающих
        циклов (например, пар вихрей, структур доменов) и потерю
        определённой \(\partial\Omega_2\), а не через спонтанное исчезновение
        упорядоченности только внутри объема.
\end{itemize}

\paragraph{Критерии фальсификации и тесты.}
\begin{itemize}
  \item \textbf{F2.1.}
        Сильный воспроизводимый фазовый переход без порогов структуры
        параметра порядка фальсифицирует P2.1.
  \item \textbf{F2.2.}
        Наблюдение устойчивых бесконечных кластеров при \(p<p_c\) в
        контролируемых перколяционно-подобных системах противоречило бы P2.2.
  \item \textbf{F2.3.}
        Демонстрация континуума, чья эффективная размерность возрастает
        (например, от 2D к 3D-поведенческим режимам), при фиксированных всех
        осей \(M_2\) и без активации новых степеней свободы, фальсифицировала
        бы P2.3 и общую теорему размерностной монотонности.
  \item \textbf{Экспериментальная программа.}\n        Точные эксперименты по разрушению когерентности, развязыванию вихрей и
        порогам перколяции могут использоваться для эмпирического картирования
        \(\Theta_{\mathrm{crit}}\) и проверки того, ведет себя
        соответствующая \(\partial\Omega_2\), как прогнозирует OC
        (замедление релаксации, локализация коллапса на границе).
\end{itemize}

% ====================================================================
\subsection{Предсказания и тесты в сценариях происхождения жизни
            (\texorpdfstring{$K_3$}{K\_3}--\texorpdfstring{$K_4$}{K\_4})}
% ====================================================================

Для сценариев происхождения жизни OC формулирует явные утверждения о
каталитическом замыкании, структуре мембраны и жизнеспособности
протоклеток.

\paragraph{Структурные предсказания.}
\begin{itemize}
  \item \textbf{P3.1 (Каталитическое замыкание).}
        Длительно устойчивые химические организации на \(K_3\) требуют как
        минимум одного цикл-комплекса, подобного RAF; системы без каталитического
        замыкания не могут поддерживать \(k_3>0\).
  \item \textbf{P3.2 (Осмотивый порог схлопывания).}
        Протоклетки с фиксированным составом мембраны подчиняются
        зависимости радиус–напряжение–проницаемость: выше критической
        комбинации осмотического градиента, кривизны и проницаемости
        схлопывание неизбежно.
  \item \textbf{P3.3 (Необходимость границы).}
        Протоклетки без структурно различимой границы
        \(\partial\Omega_4\) (мембраны или эквивалентной компартментной
        структуры) не способны поддерживать нетривиальные метаболические
        потоки \(J_{\mathrm{metabolic}}\) на масштабах, определяющих
        \(k_4>0\).
  \item \textbf{P3.4 (Сбой локализованный на патчах).}
        Схлопывание мембраны инициируется на участках границы, где локальные
        пороги первыми пересекаются, а не равномерно по всей поверхности.
\end{itemize}

\paragraph{Критерии фальсификации и тесты.}
\begin{itemize}
  \item \textbf{F3.1.}
        Демонстрация химически самоподдерживающихся организаций с
        \(k_3>0\) и без идентифицируемого каталитического замыкания или
        цикла-комплекса фальсифицировала бы P3.1.
  \item \textbf{F3.2.}
        Систематическое нарушение предсказуемых осмотических зависимостей
        схлопывания (например, произвольные градиенты, поддерживаемые без
        структурной перенастройки) противоречило бы P3.2.
  \item \textbf{F3.3.}
        Наблюдение долговечных протоклеточно-подобных систем с устойчивыми
        метаболическими потоками, но без структуры границы,
        совместимой с \(\partial\Omega_4\), фальсифицировало бы P3.3 и
        ослабило бы роль границ в OC.
  \item \textbf{Экспериментальная программа.}
        Контролируемые эксперименты протекания через протоклетки и их
        разрыва, с количественным измерением градиентов и свойств
        мембран, могут проверять структурный статус
        \(\Theta_{\mathrm{grad}}\), \(\Theta_{\mathrm{perm}}\),
        \(\Theta_{\mathrm{mem}}\) и динамики патчей.
\end{itemize}

% ====================================================================
\subsection{Предсказания и тесты в биологических системах
            (\texorpdfstring{$K_4$}{K\_4}--\texorpdfstring{$K_5$}{K\_5})}
% ====================================================================

На биологическом уровне OC фокусируется на возбудимости, протоспайках и
переходе \(K_4 \to K_5\).

\paragraph{Структурные предсказания.}
\begin{itemize}
  \item \textbf{P4.1 (Минимальный порог возбудимости).}
        Переход от чисто метаболических протоклеток к возбудимым континуумам
        требует минимального порога возбудимости
        \(\Theta_{\mathrm{exc}}\); не может существовать устойчивое
        поведение \(K_5\) без нетривиальных осей \(\Delta V\) и хорошо
        определенного порога спайка.
  \item \textbf{P4.2 (Прото-спайковые инварианты).}
        В ранних возбудимых системах существуют минимальные инварианты для
        спайкоподобных циклов (например, минимальная амплитуда
        \(\Delta V\), рефрактерный период \(\tau\)), соответствующие
        структурным ограничениям \(C_{\mathrm{spike}}\).
  \item \textbf{P4.3 (Возбудимость предшествует сложной наследственности).}
        Длительная передача сложных регуляторных структур требует предварительного
        появления каналов, подобных \(K_5\), и циклов возбудимости;
        системы, строго находящиеся на \(K_4\), не могут неограниченно
        поддерживать унаследованную регуляторную сложность без перехода на
        более высокий уровень.
\end{itemize}

\paragraph{Критерии фальсификации и тесты.}
\begin{itemize}
  \item \textbf{F4.1.}
        Демонстрация действительно \(K_5\)-подобной сигнализации (поезда
        спайков, возбудимые волны) в системах без идентифицируемых осей
        возбудимости или порогов фальсифицировала бы P4.1.
  \item \textbf{F4.2.}
        Обнаружение устойчивых сложных механизмов наследования в
        протоклеточных системах, строго ограниченных \(K_4\) (без
        возбудимости, без спайкоподобных циклов), поставило бы под
        вопрос P4.3.
  \item \textbf{Экспериментальная программа.}
        Наборы минимальных ионных каналов, синтетические возбудимые мембраны
        и контролируемые тесты нестабильности \(\Delta V\) можно использовать
        для картирования структурного пространства
        \(\Theta_{\mathrm{exc}}\), \(\Theta_{\mathrm{refr}}\) и получаемого
        \(C_{\mathrm{spike}}\).
\end{itemize}

% ====================================================================
\subsection{Предсказания и тесты в когнитивных континуумах
            (\texorpdfstring{$K_6$}{K\_6})}
% ====================================================================

Для когнитивных систем OC подчеркивает емкость связывания, ошибку предсказания
и поддержание внутренних моделей через циклы.

\paragraph{Структурные предсказания.}
\begin{itemize}
  \item \textbf{P6.1 (Лимит емкости связывания).}
        Для любого когнитивного континуума существует конечный порог емкости
        связывания \(\Theta_{\mathrm{bind}}\), выше которого дополнительное
        множество признаков нельзя интегрировать в когерентные представления
        без схлопывания \(k_6\).
  \item \textbf{P6.2 (Схлопывание по ошибке предсказания).}
        Существует порог ошибки предсказания \(\Theta_{\mathrm{pred}}\), выше
        которого не может поддерживаться когерентная внутренняя модель; когнитивное
        схлопывание наступает, когда ошибка предсказания не может быть
        снижена никакими потоками \(J_6\) при фиксированных текущих осях.
  \item \textbf{P6.3 (Схлопывание на уровне циклов).}
        Когнитивное схлопывание всегда включает разрушение ключевых циклов
        (внимания, предсказания, памяти); не существует сценария
        схлопывания при целостном \(C_{\max}(K_6)\).
\end{itemize}

\paragraph{Критерии фальсификации и тесты.}
\begin{itemize}
  \item \textbf{F6.1.}
        Надежные данные об бесконечной емкости связывания в конечных системах
        (нет наблюдаемых порогов, нет деградации) фальсифицировали бы P6.1.
  \item \textbf{F6.2.}
        Эмпирические режимы, где ошибка предсказания расходится, но устойчивые и
        когерентные внутренние модели сохраняются бесконечно, противоречили бы
        P6.2.
  \item \textbf{F6.3.}
        Демонстрация когнитивного схлопывания при сохранении целостности
        циклов внимания, предсказания и памяти фальсифицировала бы P6.3 и
        цикловое определение схлопывания.
  \item \textbf{Экспериментальная программа.}
        Поведенческие и нейрофизиологические исследования деградации
        связывания под нагрузкой, насыщения ошибки предсказания и
        реконсолидации памяти могут проверять эмпирические аналоги
        \(\Theta_{\mathrm{bind}}\), \(\Theta_{\mathrm{pred}}\) и
        \(C_{\mathrm{memory}}\).
\end{itemize}

% ====================================================================
\subsection{Предсказания и тесты в социальных/цивилизационных континуумах
            (\texorpdfstring{$K_7$}{K\_7}--\texorpdfstring{$K_8$}{K\_8})}
% ====================================================================

На более высоких социальных уровнях OC формулирует структурные ограничения на
доверие, институциональные циклы и инфраструктурную устойчивость.

\paragraph{Структурные предсказания.}
\begin{itemize}
  \item \textbf{P7.1 (Порог доверия).}
        Устойчивые социальные континуумы на \(K_7\) требуют доверия выше
        порога \(\Theta_{\mathrm{trust}}\); ниже этого порога институциональные
        циклы не могут поддерживать \(k_7>0\).
  \item \textbf{P7.2 (Замыкание институционального цикла).}
        На \(K_7\) не существует устойчивого континуума без как минимум
        одного закрытого институционального цикла (например,
        цепочки налоги\textendash услуги\textendash легитимность), совместимого
        с определением OC для \(C_7\).
  \item \textbf{P8.1 (Пороги инфраструктурных циклов).}
        Цивилизационные континуумы \(K_8\) требуют устойчивых
        инфраструктурных циклов; схлопывание на \(K_8\) включает провал границы
        в \(\partial\Omega_8\) и перколяцию отказов по ключевым
        инфраструктурным слоям.
\end{itemize}

\paragraph{Критерии фальсификации и тесты.}
\begin{itemize}
  \item \textbf{F7.1.}
        Исторические или смоделированные примеры долгоживущих институтов с
        устойчиво низким доверием относительно любого правдоподобного
        \(\Theta_{\mathrm{trust}}\) ставят под вопрос P7.1.
  \item \textbf{F7.2.}
        Демонстрация долговечных, высоко сложных обществ без идентифицируемых
        закрытых циклов управления или ресурсных потоков фальсифицировала бы
        P7.2.
  \item \textbf{F8.1.}
        Надежные модели цивилизационного схлопывания, в которых отсутствует
        разрушение инфраструктурных циклов или граница провала в
        \(\partial\Omega_8\), противоречат P8.1.
  \item \textbf{Экспериментальная программа.}
        Количественные анализы исторических данных, модели хрупкости сетей и
        крупномасштабные симуляции динамики доверия и перколяции
        инфраструктур могут использоваться для проверки структурной роли
        \(\Theta_{\mathrm{trust}}\), \(C_{\mathrm{inst}}\) и
        \(\partial\Omega_8\).
\end{itemize}

% ====================================================================
\subsection{Теоретические и метатеоретические предсказания
            (\texorpdfstring{$K_9$}{K9}--\texorpdfstring{$K_{10}$}{K10})}
% ====================================================================

OC сама и другие высокоуровневые рамки могут проверяться на теоретическом и
метатеоретическом уровнях.

\paragraph{Структурные предсказания.}
\begin{itemize}
  \item \textbf{P9.1 (Пороги когерентности).}
        Любой долговечный теоретический континуум \(K_9\) должен удовлетворять
        порогам когерентности и выразимости
        \(\Theta_{\mathrm{coh}}, \Theta_{\mathrm{expr}}\); перманентный успех
        теорий, явно нарушающих внутреннюю согласованность, противоречил бы
        рамке OC.
  \item \textbf{P9.2 (Каскады несогласованности).}
        Накопленная несогласованность в теоретической рамке должна приводить к
        схлопыванию её цикла комплекса (исследовательские программы,
        научные циклы); устойчивый прогресс при произвольно большом размере
        неконсистентности несовместим с OC.
  \item \textbf{P10.1 (Мета-неполнота).}
        Любой метатеоретический континуум \(K_{10}\) структурно неполон:
        по мере роста множеств моделей структурное напряжение на \(K_{10}\)
        должно в конечном счете превысить некоторый
        \(\Theta_{\mathrm{dim}}\), принуждая либо к расширению осей, либо
        к схлопыванию.
\end{itemize}

\paragraph{Критерии фальсификации и тесты.}
\begin{itemize}
  \item \textbf{F9.1.}
        Полностью самосогласованная закрытая теоретическая система, охватывающая
        все релевантные модели, никогда не требующая расширения или пересмотра
        и в то же время поддерживающая неограниченный эмпирический прогресс,
        ставила бы под вопрос P9.2 и P10.1.
  \item \textbf{F9.2.}
        Демонстрация теоретических рамок, которые остаются несогласованными по
        стандартам OC, но тем не менее показывают стабильный и кумулятивный
        междоменный предиктивный успех, фальсифицировала бы структурную связь
        между когерентностью и \(k_9\).
  \item \textbf{Экспериментальная программа.}
        Формальные анализы семейств моделей, логические проверки
        согласованности и метауровневые исследования смены теории могут
        использоваться для проверки существования и поведения
        \(\Theta_{\mathrm{expr}}\), \(\Theta_{\mathrm{coh}}\) и
        соответствующих циклов на уровнях \(K_9\)–\(K_{10}\).
\end{itemize}

% ====================================================================
\subsection{Мета-критерии и междоменная валидация}
% ====================================================================

В заключение OC предлагает метауровневые критерии, связывающие
фальсифицируемость между предметными областями и уровнями K.

\paragraph{Предиктивные инварианты по K-уровням.}
Предсказывается выполнение ряда структурных отношений единообразно:
\begin{itemize}
  \item существование по меньшей мере одного поддерживающего цикла
        \(C_{\mathrm{support}}\) для любого живого континуума;
  \item монотонность размерности под действием операторов рождения;
  \item необходимость совместимости пространства вложения
        \(A(K_x)\subseteq A(M_x)\) и совместимости порогов с \(M_x\);
  \item сигнатуры схлопывания: критическое замедление, локализация границы,
        разрушение циклов, расходимость структурного напряжения.
\end{itemize}
Нарушение этих инвариантов на любом хорошо изученном уровне K
распространяется вверх и вниз по иерархии.

\paragraph{Алгоритмические структурные тесты.}
Core~1.2 и последующая работа ориентированы на разработку алгоритмических
тестов для:
\begin{itemize}
  \item \(\Omega\)-согласованности (непустой допустимой области при данных
        порогах),
  \item жизнеспособности циклов (существование циклов-комплексов \(C
        ) при данных потоках),
  \item совместимости вложения (оси и пороги относительно \(M_x\)).
\end{itemize}
Эти тесты могут применяться к конкретным моделям в физике, химии,
биологии, когниции и социальных науках для проверки, являются ли они
лицензионными инстанцированиями OC.

\paragraph{Междоменная логика фальсификации.}
Поскольку аксиомы общи для всех уровней K, фальсификация имеет
\emph{междоменную} структуру:
\begin{itemize}
  \item если один структурный аксиоматический тезис опровергнут в строго
        контролируемом физическом контексте, все выводы для более высоких
        уровней должны быть пересмотрены;
  \item наоборот, согласованная валидация одного и того же структурного
        мотива (например, схлопывание через границу) в нескольких областях
        повышает доверие к соответствующей аксиоме OC.
\end{itemize}

\paragraph{Единый конвейер валидации.}
Core~1.3.3 обеспечивает структурную базу для явного конвейера валидации в
Core~1.2:
\begin{enumerate}
  \item выбираются предметные модели и выделяются их континуумные компоненты
        \((\Omega, A, P, J, \Theta, \partial\Omega, C, k)\);
  \item тестируются локальные структурные предсказания (пороги, циклы,
        границы);
  \item проверяются поперечные инварианты уровней (размерностная
        монотонность, совместимость вложения);
  \item обновляется отображение между предметными теориями и уровнями K.
\end{enumerate}

OC поэтому не является чисто концептуальной рамкой; она предназначена для
эмпирического ограничения и, в принципе, структурной фальсифицируемости по
всей иерархии \(K_0\)–\(K_{10}\).

% END OF FILE
